% SPDX-License-Identifier: CC-BY-SA-4.0
% Author: Matthieu Perrin
% Part: 
% Section: 
% Sub-section: 
% Frame: 

\begingroup

\begin{frame}[fragile]{Exercice}

  \begin{lstlisting}[numbers=left, gobble=4]
    public class WaitFreeStack<T> {
      private final AtomicInteger range = new AtomicInteger(0);
      private final AtomicReferenceArray<T> items = (*\Alert{???;}*)

      public void push(T value) {
        int i = range.getAndIncrement();
        items.set(i, value);
      }

      public T pop() {
        for(int i = range.get()-1; i>=0; i--) {
          T value = items.getAndSet(i, null);
          if(value != null) return value;
        }
        return null;
      }
    }
  \end{lstlisting}

  \begin{exampleblock}{Réécrire l'algorithme de pile wait-free en Java}
    \begin{enumerate}
    \item Remplacer le tableau infini \lstinline{items} par une liste chaînée
    \item Réduire la complexité spatiale au maximum
    \end{enumerate}
  \end{exampleblock}

\end{frame}

\endgroup
\endinput
